\par
\SourceLeaf{612}
3.\quad Der Wirkung eines geschlossenen Stroms $RR'R''R$ kann man magnetische Wirkung auf folgende Art gleich setzen. Es begrenze $RR'R''R$ eine beliebige Fläche, auf welcher man nordlichen Magnetismus nach beliebigem Verhältnisse ausgebreitet denke, mit Dichtigkeit $=\delta$. An jeder Stelle der Fläche errichte man eine unendlich kleine Normale im zusammengesetzten geraden Verhältniss der Intensität des Stromes, des verkehrten von $\delta$ und zwar nach oben oder nach unten gerichtet, je nachdem der Strom beim Umlauf um die Fläche diese rechts oder links hat. Die Endpunkte jener Normalen liegen in einer zweiten Fläche, auf welcher und in deren Theilen man genau ebenso viel südlichen Magnetismus ausbreite, als sich auf der andern und deren correspondirenden Theilen befindet. Diese zwei magnetischen Flächen aequivaliren für jeden ausser ihnen liegenden Punkt jenem galvanischen Strom.

\begin{center}[6.]\par\medskip\textit{Zur mathematischen Theorie der electrodynamischen Wirkungen.}\end{center}
1.\quad Die gegenseitige Wirkung zweier Stromelemente $\mathrm ds,\mathrm ds'$ auf einander, die Intensität der Ströme durch $i,i'$ bezeichnet, drückt \textsc{Ampère} durch die Formel
\[
ii'(\sin\theta\sin\theta'\cos\omega+k\cos\theta\cos\theta')r^n\mathrm ds\,.\,\mathrm ds'
\]
aus, indem er voraussetzt, sie habe in der verbindenden geraden Linie Statt, und positive oder negative Zeichen beziehen sich auf Anziehung oder Abstossung. Es bedeuten hier $r$ den Abstand der Elemente von einander, $\theta$ und $\theta'$ die Winkel der Elemente $\mathrm ds,\mathrm ds'$ mit $r$, letztere Linie bei beiden in gleicher Richtung genommen, endlich $\omega$ den Winkel der beiden Ebenen durch $r$ und $\mathrm ds$ einerseits, und $r$ und $\mathrm ds'$ andererseits. Aus seinen Versuchen hat \textsc{Ampère} geschlossen, dass $n=-2$, $k=-\frac{1}{2}$ gesetzt werden müsse. Die gegenseitige Anziehung wird also durch
\[
\frac{ii'(\sin\theta\sin\theta'\cos\omega-\frac{1}{2}\cos\theta\cos\theta')\mathrm ds\,.\,\mathrm ds'}{rr}
\]
gemessen, und ein negativer Werth dieses Ausdrucks bedeutet eine Abstossung.

2.\quad Wir bezeichnen noch durch $a,a'$ die ganzen Stromlängen, durch $\rho$ die Quadratwurzel aus $r$, durch $x,y,z$ die Coordinaten eines beliebigen Punkts in $s$, durch $x',y',z'$ die Coordinaten eines beliebigen Punkts in $s'$, durch $\varepsilon$ den
\SourceLeaf{613}Winkel, welchen zwei Stromelemente $\mathrm ds,\mathrm ds'$ mit einander machen. Partielle Differentiationen in Beziehung auf $s$ und $s'$ sollen durch die Charakteristiken $\mathrm d,\mathrm d'$ unterschieden werden; eben so partielle Integrationen durch die Zeichen $\int,\int'$. Wir haben
\[
\begin{gathered}
\mathrm dr=2\rho\,\mathrm d\rho=\cos\theta\,.\,\mathrm ds,\qquad \mathrm d'r=2\rho\,\mathrm d'\rho=-\cos\theta'\,.\,\mathrm ds'\\
\rho^4=(x'-x)^2+(y'-y)^2+(z'-z)^2\\
2\rho^3\mathrm d\rho=-(x'-x)\mathrm dx-(y'-y)\mathrm dy-(z'-z)\mathrm dz\\
\begin{aligned}
6\rho\rho\,\mathrm d\rho\,.\,\mathrm d'\rho+2\rho^3\mathrm d\mathrm d'\rho
&=-\mathrm dx\,.\,\mathrm d'x'-\mathrm dy\,.\,\mathrm d'y'-\mathrm dz\,.\,\mathrm d'z'\\
&=-\cos\varepsilon\,.\,\mathrm ds\,.\,\mathrm ds'\\
2\rho^3\mathrm d\mathrm d'\rho&=(-\cos\varepsilon+\tfrac{3}{2}\cos\theta\,.\,\cos\theta')\mathrm ds\,.\,\mathrm ds'\\
&=-(\sin\theta\sin\theta'\,.\,\cos\omega-\tfrac{1}{2}\cos\theta\,.\,\cos\theta')\mathrm ds\,.\,\mathrm ds'
\end{aligned}
\end{gathered}
\]
Der obige \textsc{Ampère}sche Ausdruck verwandelt sich also in
\[
-\frac{2ii'\mathrm d\mathrm d'\rho}{\rho}
\]

3.\quad Durch die Charakteristik $\delta$ bezeichnen wir die unendlich kleinen Variationen der Grösse, welcher sie vorgesetzt ist, in so fern solche von unendlich kleinen virtuellen Ortsänderungen der Punkte der Reophoren $s,s'$ abhängen. Virtuelle Ortsänderungen sind alle willkürlich gedachte, die mit den Bedingungen verträglich sind, die die Natur der Reophoren und ihre äussern Verhältnisse mit sich bringen.

4.\quad Man hat die ganze Wirkung der Ströme auf einander als bekannt anzusehen, wenn man andere Kräfte, die auf einzelne Punkte derselben in endlicher oder unendlicher Anzahl wirken, angeben kann, die ihnen aequivaliren, d.\ i.\ deren entgegengesetzte jenen das Gleichgewicht halten. Es sei $W$ das Aggregat der letztern Kräfte in ihre virtuellen Bewegungen multiplicirt. Es wird dann nach dem Princip der virtuellen Bewegungen
\[
W+\int\int'\delta r\,.\,\frac{2ii'\mathrm d\mathrm d'\rho}{\rho}=W+4ii'\int\int'\mathrm d\mathrm d'\rho\,.\,\delta\rho
\]
$=0$ sein müssen, oder wenigstens nicht positiv werden können, für jedes System von virtuellen Bewegungen der Reophoren, welches mit den Statt findenden Bedingungen verträglich ist. Für sich allein hingegen halten die Stromkräfte die Reophoren im Gleichgewicht, wenn $\int\int'\delta\rho\frac{\mathrm d\mathrm d'\rho}{\rho}$ für alle virtuelle Bewegungen $=0$ oder negativ ist. Die Integrationen sind hier von $s=0$ bis $s=a$, und von $s'=0$ bis $s'=a'$ auszuführen.

\SourceLeaf{614}5.\quad Um die Natur dieses Integrals kennen zu lernen, entwickeln wir die Variation von $\int\int'\mathrm d\rho\,.\,\mathrm d'\rho$. Wir haben
\[
\begin{aligned}
\delta\int\int'\mathrm d\rho\,.\,\mathrm d'\rho&=\int\int'\delta(\mathrm d\rho\,.\,\mathrm d'\rho)\\
&=\int\int'\delta\mathrm d\rho\,.\,\mathrm d'\rho+\int\int'\delta\mathrm d'\rho\,.\,\mathrm d\rho
\end{aligned}
\]
Nun aber ist
\[
\mathrm d\int'\mathrm d'\rho\delta\rho=\int'\mathrm d\delta\rho\,.\,\mathrm d'\rho+\int'\delta\rho\,.\,\mathrm d\mathrm d'\rho
\]
oder wenn man von $s=0$ bis $s=a$ integrirt und die Werthe von $\int'\mathrm d'\rho\,.\,\delta\rho$ für $s=0$ und $s=a$ durch $F',G'$,
\[
-F'+G'=\int\int'\mathrm d\delta\rho\,.\,\mathrm d'\rho+\int\int'\delta\rho\,.\,\mathrm d\mathrm d'\rho
\]
und eben so wenn man die Werthe von $\int\mathrm d\rho\,.\,\delta\rho$ für $s'=0$, und $s'=a'$ durch $F,G$ bezeichnet
\[
-F+G=\int\int'\mathrm d'\delta\rho\,.\,\mathrm d\rho+\int\int'\delta\rho\,.\,\mathrm d\mathrm d'\rho
\]
Folglich
\[
\begin{aligned}
-F-F'+G+G'&=\int\int'\mathrm d\delta\rho\,.\,\mathrm d'\rho+\int\int'\mathrm d'\delta\rho\,.\,\mathrm d\rho\\
&\qquad+2\int\int'\delta\rho\,.\,\mathrm d\mathrm d'\rho\\
&=\delta\int\int'\mathrm d\rho\,.\,\mathrm d'\rho+2\int\int'\delta\rho\,.\,\mathrm d\mathrm d'\rho
\end{aligned}
\]

Es ist also das virtuelle Moment der gegenseitigen Einwirkung der Stromstücke $s,s'$ auf einander
\[
\begin{aligned}
V=4ii'\int\int'\mathrm d\mathrm d'\rho\,.\,\delta\rho&=-2ii'\delta\int\int'\mathrm d\rho\,.\,\mathrm d'\rho\\
&\qquad+2ii'(G+G'-F-F')
\end{aligned}
\]

6.\quad Eine stromerzeugende oder electromotorische Kraft übt ein Strom nur aus, indem er entsteht oder indem er sich bewegt. Die electromotorische Kraft eines Stromelements $i\,\mathrm ds$ wirkt in jedem Punkt mit einer Stärke, welche der Entfernung $r$ verkehrt, hingegen dem auf diese Linie projicirten Stromelement direct proportional ist, und in der Richtung der Linie $r$ selbst, aber stets im entgegengesetzten Sinn. Man schliesst hieraus leicht, dass das virtuelle Moment der electromotorischen Kraft des Elements $i\,\mathrm ds$ durch
\[
+\frac{i\,\mathrm ds}{r}\,.\,\frac{\mathrm dr}{\mathrm ds}\,.\,\delta r
\]
ausgedrückt werden kann oder durch
\SourceLeaf{615}\[
4i\,\mathrm d\rho\,.\,\delta\rho
\]

In dem Rheophorelement $\mathrm ds'$ ist also die electromotorische Kraft des Stromelements $i\,\mathrm ds$
\[
=4i\,\mathrm d\rho\,.\,\frac{\mathrm d\rho}{\mathrm ds'}\,.\,\mathrm ds'=4i\,\mathrm d\rho\,.\,\mathrm d'\rho
\]

Das doppelte Integral $4i\int\int'\mathrm d\rho\,.\,\mathrm d'\rho$ ist also die ganze electromotorische Kraft in dem Rheophorstück $s'=0$ bis $s'=a'$, welche durch das Stromstück $is$, von $s=0$ bis $s=a$ ausgeübt wird. Wir bezeichnen diese mit $A$. Ferner ist das virtuelle Moment der electromotorischen Kraft, welche dasselbe Stromstück in einem gegebenen Punkte ausübt
\[
=4i\int\mathrm d\rho\,.\,\delta\rho\qquad\text{von }s=0\text{ bis }s=a
\]
ist diese, in dem Anfang des Rheophorstücks $s'$, $=B$, an dessen Ende $=C$, so ist
\[
B=4iF,\qquad C=4iG
\]
Bedeuten ebenso $B',C'$ die virtuellen Momente der electromotorischen Kraft, welche ein Stromstück $i's'$ von $s'=0$ bis $s'=a'$ in den Punkten Anfang und Ende von $s$ ausübt, so ist $B'=4i'F'$, $C'=4i'G'$.

Wir haben folglich
\[
V=-\tfrac{1}{2}i'\delta A-\tfrac{1}{2}i'B+\tfrac{1}{2}i'C-\tfrac{1}{2}iB'+\tfrac{1}{2}iC'
\]

\begin{center}[7.]\par\medskip\textit{Das Inductionsgesetz.}\par\medskip
{\small (Gefunden 1835. Januar 23. Morgens $7^{\mathrm u}$ v.\ d.\ Aufst.)}
\end{center}
I.\quad Die Stromerzeugende Kraft, welche in einem Punkte $P$ hervorgebracht wird durch ein Rheophorelement $\gamma$, dessen Entfernung von $P$, $=r$, ist während des Zeitelements $\mathrm dt$ die Differenz der beiden Werthe von $\frac{\gamma}{r}$, welche den Augenblicken $t$ und $t+\mathrm dt$ entsprechen, durch $\mathrm dt$ dividirt, wo $\gamma$ nach Grösse und Richtung zu berücksichtigen ist, was kurz und verständlich durch
\[
-\frac{\mathrm d\dfrac{\gamma}{r}}{\mathrm dt}
\]
ausgedrückt werden kann.

\SourceLeaf{616}II.\quad Die Stärke eines erzeugten Stroms ist
\[
=\frac{\int p\,\mathrm ds}{\theta}
\]
wo $p$ die Stromerzeugende Kraft in jedem Element $\mathrm ds$ des Rheophors, $\theta$ der ganze Widerstand.

\begin{center}[8.]\end{center}
Es seien $s,s'$ zwei geschlossene Rheophoren, $r$ die gegenseitige Distanz zweier Punkte in $s$ und $s'$, $r=\rho\rho$, $\theta$ der ganze Widerstand in $s'$.

Folgendes sind die beiden Grundgesetze:

I.\quad Sind in $s$ und $s'$ galvanische Ströme mit den Intensitäten $\varepsilon,\varepsilon'$, die als positiv betrachtet werden, wenn sie die Rheophoren in dem Sinn durchlaufen, in welchem deren Elemente $\mathrm ds,\mathrm ds'$ gewählt werden, so ist die abstossende Kraft der Elemente
\[
=+\frac{\varepsilon\varepsilon'}{\rho}\,.\,\frac{\mathrm d\mathrm d\rho}{\mathrm ds\,.\,\mathrm ds'}\,.\,\mathrm ds\,.\,\mathrm ds'
\]
oder wenn man durch $\mathrm d,\mathrm d'$ die partiellen auf beide Ströme sich beziehenden Differentiationen bezeichnet
\[
=+\frac{\varepsilon\varepsilon'}{\rho}\,.\,\mathrm d\mathrm d'\rho
\]
Diese Kraft wirkt in der Richtung der geraden Linie $r$.

\vspace{3mm}
II.\quad Entsteht während der sehr kleinen Zeit $\delta t$ der Strom in $s$, so ist damit eine oben bemerkte Stromerzeugende Kraft in jedem Punkte begleitet; vom Element $\mathrm ds'$ ist das Maass derselben
\[
=-\frac{\varepsilon\,\mathrm ds\,.\,\mathrm ds'}{\delta t\,.\,r}\cos u
\]
wenn $u$ die Neigung der Richtungen $\mathrm ds,\mathrm ds'$ gegen einander bezeichnet.

Bezeichnet man durch $z$ die Projection von $r$ auf die Richtung von $\mathrm ds'$, so ist $\mathrm ds\,.\,\cos u=\mathrm dz$, also jene Formel
\[
=-\frac{\varepsilon\,\mathrm dz}{\delta t\,.\,r}\,.\,\mathrm ds'
\]
oder die ganze aus $is$ in $\mathrm ds'$ erregte Stromerzeugende Kraft
\SourceLeaf{617}\[
=-\frac{\varepsilon\,.\,\mathrm ds'}{\delta t}\int\frac{\mathrm dz}{r}
\]

Da $\frac{\mathrm dz}{r}-\frac{z\,\mathrm dr}{rr}$ ein vollständiges Differential, mithin dessen Integral durch den geschlossenen Rheophor $s$ ausgedehnt $=0$ ist, so ist obiges Integral auch
\[
=-\frac{\varepsilon\,\mathrm ds'}{\delta t}\int\frac{z\,\mathrm dr}{rr}
\]
oder da $-\frac{z}{r}\,.\,\mathrm ds'=\mathrm d'r$
\[
=+\frac{\varepsilon}{\delta t}\int\frac{\mathrm dr\,.\,\mathrm d'r}{r}=+\frac{4\varepsilon}{\delta t}\int\mathrm d\rho\,.\,\mathrm d'\rho
\]
oder die ganze Stromerzeugende Kraft $=\frac{4\varepsilon}{\delta t}\int\int'\mathrm d\rho\,.\,\mathrm d'\rho$ folglich hat der in $s'$ während der Zeit $\delta t$ Statt findende Strom die Intensität
\[
\frac{4\varepsilon}{\theta\delta t}\,.\,\int\int'\mathrm d\rho\,.\,\mathrm d'\rho
\]
wofür man auch offenbar
\[
-\frac{4\varepsilon}{\theta\delta t}\,.\,\int\int'\rho\,\mathrm d\mathrm d'\rho
\]
schreiben kann. Es ist nemlich $\int\,.\,\int'\mathrm d\rho\,.\,\mathrm d'\rho=\int\rho\,\mathrm d\rho-\int\int'\rho\,\mathrm d\mathrm d'\rho$ und $\int\rho\,\mathrm d\rho=0$ indem es durch die ganze Stromlinie $s$ ausgedehnt wird.

\begin{center}[9.]\par\medskip\textit{Induction durch Bewegung.}\end{center}
I.\quad In jedem Punkte des Raumes, dessen Coordinaten $x,y,z$, bezeichne $V$ den körperlichen Winkel, welchen ein Strom $S'$ in jenem Punkte umspannt.

$V=\text{Const.}$ bestimmt daher eine Fläche, deren tangirende Ebene und Normale dasselbe sind, was \textsc{Ampère} Plan directeur. und Directrice nennt.

II.\quad In jedem bewegten körperlichen Molecule $\mu$, dessen partielle Geschwin\-digkeiten $\frac{\mathrm dx}{\mathrm dt},\frac{\mathrm dy}{\mathrm dt},\frac{\mathrm dz}{\mathrm dt}$ sind, findet eine electrodynamische \textit{Kraft} Statt, deren partielle Zerlegungen $\xi,\eta,\zeta$ sein mögen. Man hat dann
\[
\begin{aligned}
\xi&=\left(\frac{\mathrm dV}{\mathrm dy}\,.\,\frac{\mathrm dz}{\mathrm dt}-\frac{\mathrm dV}{\mathrm dz}\,.\,\frac{\mathrm dy}{\mathrm dt}\right)\mu\\[4pt]
\eta&=\left(\frac{\mathrm dV}{\mathrm dz}\,.\,\frac{\mathrm dx}{\mathrm dt}-\frac{\mathrm dV}{\mathrm dx}\,.\,\frac{\mathrm dz}{\mathrm dt}\right)\mu\\[4pt]
\zeta&=\left(\frac{\mathrm dV}{\mathrm dx}\,.\,\frac{\mathrm dy}{\mathrm dt}-\frac{\mathrm dV}{\mathrm dy}\,.\,\frac{\mathrm dx}{\mathrm dt}\right)\mu
\end{aligned}
\]

\SourceLeaf{618}III.\quad Geht hingegen durch jenen Punkt ein Element eines Stromkörpers $\mathrm ds$, so wird solcher von $S$ sollicitirt, und sind die partiellen Kräfte $\xi,\eta,\zeta$, so ist, $i$ Intensität des Stroms,
\[
\begin{aligned}
\xi&=\left(\frac{\mathrm dV}{\mathrm dy}\,.\,\frac{\mathrm dz}{\mathrm ds}-\frac{\mathrm dV}{\mathrm dz}\,.\,\frac{\mathrm dy}{\mathrm ds}\right)i\,\mathrm ds\\
&\text{etc.}
\end{aligned}
\]

Es ist übrigens
\[
\begin{aligned}
V&=\int\frac{z'(x'\mathrm dy'-y'\mathrm dx')}{r'(x'x'+y'y')}=\int\frac{x'(y'\mathrm dz'-z'\mathrm dy')}{r'(y'y'+z'z')}\\
&=\int\frac{y'(z'\mathrm dx'-x'\mathrm dz')}{r'(z'z'+x'x')}
\end{aligned}
\]
wenn $x',y',z'$ sich auf den \textit{wirkenden} Strom $S'$ beziehen, woraus
\[
\frac{\mathrm dV}{\mathrm dz}=\int\frac{x'\mathrm dy'-y'\mathrm dx'}{r'^3}
\]
welches \textsc{Ampères} Ausdruck ist. Es ist indefinit
\[
\int\left\{\frac{z(x\mathrm dy-y\mathrm dx)}{r(xx+yy)}-\frac{y(z\mathrm dx-x\mathrm dz)}{r(xx+zz)}\right\}=\text{Arc.}\,\operatorname{tg}\frac{yz}{xr}
\]
am einfachsten bewiesen indem man es in die Form setzt
\[
\frac{-yzrr\,\mathrm dx-xyzx\,\mathrm dx+xzrr\,\mathrm dy-xyzy\,\mathrm dy+xyrr\,\mathrm dz-xyzz\,\mathrm dz}{r(xx+yy)(xx+zz)}
\]

\vspace{4mm}
Die ganze, von der \textit{Entstehung eines Stroms} herrührende Stromerzeugende Kraft in jedem Punkte des Raums, z.\ B.\ in dem Punkte, dessen Coordinaten alle $=0$, wird in drei partielle Kräfte zerlegt, nemlich
\[
\begin{aligned}
X&=\int\frac{\mathrm dx'}{r}&\quad\text{also }\frac{\mathrm dV}{\mathrm dx}&=\frac{\mathrm dZ}{\mathrm dy}-\frac{\mathrm dY}{\mathrm dz}\\[4pt]
Y&=\int\frac{\mathrm dy'}{r}&\frac{\mathrm dV}{\mathrm dy}&=\frac{\mathrm dX}{\mathrm dz}-\frac{\mathrm dZ}{\mathrm dx}\\[4pt]
Z&=\int\frac{\mathrm dz'}{r}&\frac{\mathrm dV}{\mathrm dz}&=\frac{\mathrm dY}{\mathrm dx}-\frac{\mathrm dX}{\mathrm dy}
\end{aligned}
\]

(Februar 4.) [Späterer Zusatz:] (1836 April 7.)

(Für plane Curve ist $z'$ constant also $Z=0$. Es sei
\[
\int r\frac{(x'-x)\mathrm dy'-(y'-y)\mathrm dx'}{(x'-x)^2+(y'-y)^2}=U
\]
\[
U=\int r\,\mathrm d'\lambda=(z'-z)\int\frac{\mathrm d'\lambda}{\sin\beta}\quad\text{wenn}\quad z'-z=r\sin\beta,\quad\frac{y'-y}{x'-x}=\operatorname{tg}\lambda
\]
\SourceLeaf{619}dann erhält man
\[
\begin{aligned}
X&=\frac{\mathrm dU}{\mathrm dy}\\[4pt]
Y&=-\frac{\mathrm dU}{\mathrm dx}
\end{aligned}
\]

Die Richtung der Kraft, deren Componenten $X,Y,Z$, fällt also immer in die Fläche, wofür $U=\text{const.}$ und zugleich in das Planum, wofür $z$ constant. Diese Linie kehrt also in sich selbst zurück.)

Man hat
\[
\begin{aligned}
\xi&=\frac{\mathrm dV}{\mathrm dy}\,.\,\mathrm dz-\frac{\mathrm dV}{\mathrm dz}\,.\,\mathrm dy\\[4pt]
&=\int'\frac{\begin{gathered}((z'-z)\mathrm dx'-(x'-x)\mathrm dz')\mathrm dz\\-((x'-x)\mathrm dy'-(y'-y)\mathrm dx')\mathrm dy\end{gathered}}{r^3}\\[4pt]
&=\int'\left\{\begin{gathered}\frac{\mathrm dx'[(x'-x)\mathrm dx+(y'-y)\mathrm dy+(z'-z)\mathrm dz]}{r^3}\\-\frac{(x'-x)(\mathrm dx\,\mathrm dx'+\mathrm dy\,\mathrm dy'+\mathrm dz\,\mathrm dz')}{r^3}\end{gathered}\right\}
\end{aligned}
\]
also
\[
\begin{aligned}
\xi\delta x+\eta\delta y+\zeta\delta z&=-\int'\frac{\mathrm dr}{rr}(\mathrm dx'\,.\,\delta x+\mathrm dy'\,.\,\delta y+\mathrm dz'\,.\,\delta z)\\
&\qquad+\int'\frac{\delta r}{rr}(\mathrm dx'\,.\,\mathrm dx+\mathrm dy'\,.\,\mathrm dy+\mathrm dz'\,.\,\mathrm dz)\\
&=-\int'\delta\,.\,\frac{\mathrm dx\,.\,\mathrm dx'+\mathrm dy\,.\,\mathrm dy'+\mathrm dz\,.\,\mathrm dz'}{r}\\
&\qquad+\int'\left(\mathrm dx'\,.\,\frac{\delta\mathrm dx}{r}+\mathrm dy'\,.\,\frac{\delta\mathrm dy}{r}+\mathrm dz'\,.\,\frac{\delta\mathrm dz}{r}\right)\\
&\qquad-\int'\left(\mathrm dx'\,.\,\frac{\delta x\,.\,\mathrm dr}{rr}+\mathrm dy'\,.\,\frac{\delta y\,.\,\mathrm dr}{rr}+\mathrm dz'\,.\,\frac{\delta z\,.\,\mathrm dr}{rr}\right)
\end{aligned}
\]
die beiden letzten Theile werden
\[
\int'\mathrm dx'\,.\,\mathrm d\frac{\delta x}{r}+\mathrm dy'\,.\,\mathrm d\frac{\delta y}{r}+\mathrm dz'\,.\,\mathrm d\frac{\delta z}{r}
\]
welches durch ganz $S$ integrirt verschwindet. Man hat also
\[
\int\xi\delta x+\eta\delta y+\zeta\,\mathrm dz=-\delta\int\int'\frac{\mathrm dx\,.\,\mathrm dx'+\mathrm dy\,.\,\mathrm dy'+\mathrm dz\,.\,\mathrm dz'}{r}
\]

Kürzer wird der Beweis so geführt. Man hat
\[
\begin{aligned}
\text{1)}\qquad\delta\int\int\rho\,\mathrm d\mathrm d'\rho&=\int\int\delta\rho\,.\,\mathrm d\mathrm d'\rho+\int\int\rho\delta\mathrm d\mathrm d'\rho\\[4pt]
\text{2)}\qquad\int\int\rho\,\mathrm d'\mathrm d\delta\rho+\int\int\mathrm d'\rho\,.\,\mathrm d\delta\rho&=\int\int\mathrm d'(\rho\,\mathrm d\delta\rho)=0\\[4pt]
\text{3)}\qquad0=\int\int\mathrm d(\delta\rho\,.\,\mathrm d'\rho)&=\int\int\mathrm d'\rho\,.\,\mathrm d\delta\rho+\int\int\mathrm d\mathrm d'\rho\,.\,\delta\rho
\end{aligned}
\]
\SourceLeaf{620}also durch Addition
\[
\delta\int\int\rho\,\mathrm d\mathrm d'\rho=2\int\int\delta\rho\,.\,\mathrm d\mathrm d'\rho=\int\int\delta r\,.\,\frac{\mathrm d\mathrm d'\rho}{\rho}
\]

Das unter dem Variationszeichen stehende Doppel-Integral kann auf verschiedene Arten ausgedrückt werden
\[
\begin{aligned}
\int\int\rho\,\mathrm d\mathrm d'\rho&=-\int\int\mathrm d\rho\,.\,\mathrm d'\rho\\
&=\int\int\frac{1}{4r}\cos\theta\,.\,\cos\theta'\,.\,\mathrm ds\,.\,\mathrm ds'\\
&=\int\int\frac{1}{4r}(\cos\theta\,.\,\cos\theta'+\sin\theta\sin\theta'\cos\omega)\mathrm ds\,.\,\mathrm ds'\\
&=\frac{1}{m}\,.\,\frac{1}{2-m}\int\int\rho^{2-m}\mathrm d\mathrm d'\rho^m\\
&=-\frac{1}{mm}\int\int\rho^{2-2m}\mathrm d\rho^m\,.\,\mathrm d'\rho^m
\end{aligned}
\]

\begin{center}[10.]\par\medskip\textit{Einfachste Ausdrücke für die Wirkungen galvanischer Ströme.}\end{center}
Die Fundamentalebene geht durch das wirkende Stromelement $AB$ und den Punkt auf welchen gewirkt wird $C$.

Die complexen Grössen, welche die Plätze $B,C$ relativ gegen $A$ bezeichnen, seien $\beta,\gamma$; ferner sei $r$ der Modul von $\beta$. Endlich, falls auch in $C$ ein Strom in dessen Element $CD$ bereits vorhanden, sei $\gamma+\delta+i'\zeta$ die complexe Grösse, die den Platz von $D$ gegen $A$ bezeichnet. Man hat dann

I.\quad Wenn in $C$ ein Strom ist, für die Kraft, welche dessen materieller Träger durch $AB$ erleidet
\[
\frac{\delta}{r}\quad\mathrm{Im.}\quad\frac{\beta}{\gamma}
\]

II.\quad Wenn in $C$ kein Strom ist, aber eine Bewegung in $C$ Statt findet die durch $\varepsilon$ in dem durch $BC$ gehenden Planum bezeichnet wird, die electromotorische Kraft
\[
\frac{\beta}{r}\quad\mathrm R\quad\frac{\varepsilon}{\gamma}
\]
oder in übersichtlichern Zeichen
\SourceLeaf{621}\[
\begin{gathered}
G\quad\text{wirkendes Stromelement}\\[4pt]
\left.\begin{gathered}g\\m\end{gathered}\right\}\quad
\left.\begin{gathered}\text{vorhandenes Stromelement}\\\text{vorhandene Bewegung}\end{gathered}\right\}\quad
\begin{gathered}\text{in dem Punkte auf}\\\text{welchen gewirkt wird}\end{gathered}\\[4pt]
\left.\begin{gathered}\gamma\\\mu\end{gathered}\right\}\quad
\text{Kraft zur Erregung}\Bigg\}\quad
\begin{gathered}\text{eines Stroms}\\\text{einer Bewegung}\end{gathered}\\[4pt]
r\quad\text{Entfernung als Modul der complexen Grösse }l
\end{gathered}
\]
\[
\begin{aligned}
\text{1)}\qquad r\mu&=g\,.\,\mathrm{J.}\,\frac{G}{l}\\[4pt]
\text{2)}\qquad r\gamma&=G\,.\,\mathrm{R.}\,\frac{m}{l}
\end{aligned}
\]

\begin{center}[11.]\par\medskip\textit{Geradlinige Polygone.}\end{center}
Der Punkt auf welchen gewirkt wird sei der Nullpunkt, dann ist
\[
\text{I. für}\qquad X=\int\frac{\mathrm dx}{r}
\]
der Betrag aus der ersten Seite $PP'$
\[
\frac{x'-x}{r}\log\frac{\operatorname{cotg}\tfrac{1}{2}\theta'}{\operatorname{cotg}\tfrac{1}{2}\theta}
\]
wenn $\theta,\theta'$ die Winkel zwischen $PP'$ und $0P,0P'$ sind, $PP'$ in gleicher Richtung verstanden. Die Grösse unter der Characteristik log.\ ist
\[
\begin{aligned}
&=\frac{r'+\dfrac{(x'-x)x'+(y'-y)y'+(z'-z)z'}{PP'}}{r+\dfrac{(x'-x)x+(y'-y)y+(z'-z)z}{PP'}}\\[6pt]
&=\frac{PP'\,.\,r'+r'r'-(r',PP')}{PP'\,.\,r-rr+(r,PP')}
\end{aligned}
\]

II.\quad Für $V$ der Betrag aus dem Winkel an $P'$, der Unterschied des Winkels zwischen den Ebenen $0PP',0P'P''$ von $180^0$. Der Flächeninhalt eines sphärischen Dreiecks, dessen Seiten $a,b,c$, $=\omega$ gesetzt ist
\[
\sin\tfrac{1}{2}\omega=\frac{6\text{ Pyramide }0PP'P''}{4rr'r'\cos\tfrac{1}{2}a\,.\,\cos\tfrac{1}{2}b\,.\,\cos\tfrac{1}{2}c}
\]
